% \section{The Case for Private Agents}
\section{Towards Private Agents in Multi-Agent Collaboration}
\label{sec:problem_case}

\begin{comment}
-> 반토막 내버리기
As AI systems become increasingly personalized, we are entering an era where each individual or group owned a AI agent~\citep{li2024personalLLMagents, richardson2023integrating, salemi2024optimization, tan2024democratizing}. 
However, unlike traditional multi-agent collaboration~\citep{selforganized, fourney2024magenticone, qian2024scaling, he2024llmbasedmultiagent}, private agents are fundamentally constrained by privacy. 
These agents act on behalf of their users—managing information, making recommendations, negotiating resources, and coordinating actions with other agents. 
In this work, we refer to such agents as \emph{private agents}: agents that serve a specific owner. 
In this situation, collaboration among private agents becomes inevitable, particularly when tasks span multiple owners or distributed sources of information.
As shown in Figure~\ref{fig:private_agents_collaboration}, each private agent maintains information that is sensitive, personal, or proprietary, and therefore cannot be freely shared even in situations where information exchange might improve collective performance. 
As a result, collaboration must take place under persistent information asymmetry with privacy constraints, where agents pursue shared objectives.



Our goal is to ... 지랄지랄
Importantly, these privacy constraints are intrinsic characteristics rather than incidental limitations because each agent is designed to serve and protect the interests of a specific owner.
그러기 위해서는 we believe that 
Consequently, collaboration of the private agents cannot rely on full observability, as is often assumed in classical multi-agent collaboration.
Despite these constraints, agents are still expected to contribute meaningfully to shared collaborative goals such as jointly drafting documents, constructing shared databases and common scheduling.
As shown in Figure~\ref{fig:private_agents_collaboration}, agents must decide how to participate and contribute while preserving control over the information that remains internal.
The central challenge, therefore, is not whether collaboration is possible, but how collaboration unfolds when privacy is treated as a constraint. 
Privacy-aware collaboration requires agents to balance two competing demands: advancing a common objective and respecting privacy constraints. 
Understanding this balance is essential for evaluating collaborative agent systems intended for real-world deployment, where privacy preservation is not optional but fundamental.
\end{comment}

\begin{comment}
    

As AI systems become increasingly personalized, we are entering an era where each individual or group owned a AI agent~\citep{li2024personalLLMagents, richardson2023integrating, salemi2024optimization, tan2024democratizing}. 
These agents act on behalf of their users—managing information, making recommendations, negotiating resources, and coordinating actions with other agents. 
In this work, we refer to such agents as \emph{private agents}: agents that serve a specific owner. 
In this situation, collaboration among private agents becomes inevitable, particularly when tasks span multiple owners or distributed sources of information.

However, unlike traditional multi-agent collaboration~\citep{selforganized, fourney2024magenticone, qian2024scaling, he2024llmbasedmultiagent}, private agents are fundamentally constrained by privacy. 
As shown in Figure~\ref{fig:private_agents_collaboration}, each private agent maintains information that is sensitive, personal, or proprietary, and therefore cannot be freely shared even in situations where information exchange might improve collective performance. 
As a result, collaboration must take place under persistent information asymmetry with privacy constraints, where agents pursue shared objectives.

Importantly, these privacy constraints are intrinsic characteristics rather than incidental limitations because each agent is designed to serve and protect the interests of a specific owner.
Consequently, collaboration of the private agents cannot rely on full observability, as is often assumed in classical multi-agent collaboration.
Despite these constraints, agents are still expected to contribute meaningfully to shared collaborative goals such as jointly drafting documents, constructing shared databases and common scheduling.
As shown in Figure~\ref{fig:private_agents_collaboration}, agents must decide how to participate and contribute while preserving control over the information that remains internal.

The central challenge, therefore, is not whether collaboration is possible, but how collaboration unfolds when privacy is treated as a constraint. 
Privacy-aware collaboration requires agents to balance two competing demands: advancing a common objective and respecting privacy constraints. 
Understanding this balance is essential for evaluating collaborative agent systems intended for real-world deployment, where privacy preservation is not optional but fundamental.
\end{comment}

% 미래를 봤을때의 필요성을 강조
\begin{comment}
AI agent의 Personalization 에 대한 수요가 증가, 미래를 봤을때 곧 individual 이나 organization 이 own agent가 있을것임.
conventional 한 multi-agent system은 협력, 토론 등으로 공동의 문제를 수행하는데에 집중했으나, individual 이나 organization 이 own agent를 가지고 있다고 생각했을때 가장 큰 차이는 privacy 정보라고 생각함.
For example, as shown in Figure~\ref{}, Company A 는 최선의 거래를 하기 위해 E-commerce 회사 중 하나를 선택하고 싶은데, Company C 와의 schedule 이 있는것을 Company B 가 아는것을 원치 않을 수 있음. 
하지만, 현재 agent들은 privacy 한 내용도 말할수있음.
따라서 우리는 agents that serve a specific owner, managing information, making recommendations, and coordinating actions on their behalf 하지만 privacy constraint 은 지키는 agent를 refer as \emph{private agents.}


As the demand for personalized AI systems grows~\citep{li2024personalLLMagents, richardson2023integrating, salemi2024optimization, tan2024democratizing}, we are approaching a future where every individual and organization will have their own AI agent.

While conventional multi-agent systems have focused on collaboration and debate to solve shared problems~\citep{selforganized, fourney2024magenticone, qian2024scaling, he2024llmbasedmultiagent}, the key difference when individuals and organizations own their agents lies in privacy. 
For example, as shown in Figure~\ref{fig:private_agents_collaboration}, Company A may want to select the best e-commerce partner for a transaction, but does not want Company B to know about its existing schedule with Company C. However, current agents may inadvertently reveal such private information during collaboration. We therefore introduce \emph{private agents}—agents that serve a specific owner, managing information, making recommendations, and coordinating actions on their behalf, while respecting privacy constraints that limit what information can be disclosed.

\end{comment}

With AI agents increasingly tailored to individual needs and organizational contexts~\citep{richardson2023integrating, li2024personalLLMagents, salemi2024optimization, kwon2025embodied}, we are moving toward a new era where every individual and organization has their own AI agent that communicates with other users or AI agents. 
This vision of individually-owned agents differs fundamentally from current multi-agent system approaches, where the focus is solely on successful collaboration~\citep{selforganized, fourney2024magenticone, qian2024scaling, he2024llmbasedmultiagent}, without considering privacy concerns related to personal information or proprietary data.
For instance, as shown in Figure~\ref{fig:private_agents_collaboration}, a logistics company may seek to identify the most suitable e-commerce partner for a transaction, while being unable to disclose existing contractual commitments with another company.
Here, effective collaboration requires agents to negotiate, infer, and decide with incomplete and selectively revealed information.
Such scenarios reveal that existing multi-agent approaches may not adequately handle the privacy requirements inherent in real-world collaborations.


This motivates our focus on \emph{private agents}, which we define as agents that serve individual owners, managing their information and coordinating actions while maintaining privacy constraints.
To build private agents that can operate in complex real-world settings, we must first examine how current agents perform—and where they fall short—in multi-agent collaboration with privacy constraints.
Therefore, we provide a reliable benchmark that can serve as a foundation for developing private agents, enabling future progress in this direction.

\begin{comment}
    

With AI agents increasingly tailored to individual needs and organizational contexts~\citep{richardson2023integrating, li2024personalLLMagents, salemi2024optimization, tan2024democratizing, kwon2025embodied}, we are moving toward a new era where every individual and organization has their own AI agent that communicates with other users or AI agents. 
Consequently, collaboration takes place among independently owned agents, each operating with different memory and privacy constraints.

This vision of individually-owned agents differs fundamentally from current multi-agent system approaches, where the focus is solely on successful collaboration~\citep{selforganized, fourney2024magenticone, qian2024scaling, he2024llmbasedmultiagent}, without considering privacy concerns related to personal information or proprietary data.
As a result, strong collaborative behavior in these environments often relies on assumptions that do not hold when agents act on behalf of real individuals or organizations.

In realistic collaborations, agents must reason under strict privacy constraints. 
As illustrated in Figure~\ref{fig:private_agents_collaboration}, a logistics company may seek to identify the most suitable e-commerce partner for a transaction, while being unable to disclose existing contractual commitments with another company. 
Here, effective collaboration requires agents to negotiate, infer, and decide with incomplete and selectively revealed information. Importantly, performance cannot be measured solely by whether collaboration succeeds, but also by how much sensitive information is exposed in the process.

\graytext{
These considerations reveal a missing evaluation axis in current multi-agent research: the ability to collaborate effectively despite privacy constraints that are intrinsic rather than optional. Without benchmarks that explicitly model private ownership, asymmetric information, and disclosure trade-offs, it remains unclear how well existing agent architectures would perform in such environments. This gap motivates the need for benchmark settings that treat privacy-aware collaboration not as an auxiliary concern, but as a first-class challenge.
}
\end{comment}



% \graytext{
% We refer to these as \emph{private agents}—agents that serve a specific owner, managing information, making recommendations, and coordinating actions on their behalf. 
% private agents are fundamentally constrained by privacy: they maintain information that is sensitive, personal, or proprietary, and cannot be freely shared even when disclosure might improve collective outcomes. Yet collaboration among private agents becomes inevitable when tasks span multiple owners or require distributed sources of information. As illustrated in Figure~\ref{fig:private_agents_collaboration}, this creates a challenging scenario where agents must pursue shared objectives under persistent information asymmetry, balancing task success with privacy.
% }

\begin{comment}
% 실제로는 어떻게 되어야하냐에 대한 우리의 견해
\graytext{
Private agents face a fundamental tension: they must collaborate toward shared goals while maintaining privacy constraints that are intrinsic to their design—each agent is built to serve and protect the interests of a specific owner. This means that collaboration cannot rely on full observability, as is often assumed in classical multi-agent systems~\citep{}. Yet despite these constraints, agents are still expected to contribute meaningfully to collaborative tasks such as jointly drafting documents, constructing shared databases, or coordinating schedules. As shown in Figure~\ref{fig:private_agents_collaboration}, agents must decide how much to reveal and how to participate while preserving control over sensitive information. The central challenge is therefore not whether collaboration is possible, but how it unfolds when privacy is treated as a hard constraint. Understanding this balance—between advancing common objectives and respecting privacy—is essential for evaluating agent systems intended for real-world deployment, where privacy preservation is not optional but fundamental.
}
\end{comment}